\fichatitulo{09 · FUNDAMENTO DE RAG}{Artigo de conferência · 2020}{RAG: formulação original}
{\small LEWIS, Patrick et al. \textit{Retrieval-Augmented Generation for Knowledge-Intensive NLP Tasks}. Advances in Neural Information Processing Systems, v. 33, 2020, p. 9459--9474. \href{https://proceedings.neurips.cc/paper/2020/hash/6b493230205f780e1bc26945df7481e5-Abstract.html}{Fonte e versão}.\par}

\campo{Problema e objetivo}
Como combinar o conhecimento armazenado nos parâmetros do modelo com uma memória documental recuperável para tarefas intensivas em conhecimento?

\campo{Principal contribuição · síntese interpretativa}
Formaliza RAG-Sequence e RAG-Token, que marginalizam documentos recuperados na geração. Integra recuperação neural com geração seq2seq e ajusta os componentes no aprendizado (seção 2).

\campo{Método}
Usa índice da Wikipédia de dezembro de 2018 e tarefas de perguntas e respostas, geração de perguntas e verificação de afirmações. Compara modelos paramétricos e modelos com recuperação (seções 3--4).

\campo{Resultado principal}
Em 452 pares de perguntas geradas para Jeopardy, avaliadores consideram RAG mais factual em 42,7\% dos casos, BART em 7,1\% e ambos factuais em 17\%. São categorias específicas da comparação, não uma taxa geral de respostas corretas (seção 4.3).

\campo{Excertos-chave · seleção literal}
\begin{itemize}
\excerto{Memória}{parametric and non-parametric memory}{Resumo.}
\excerto{Dependência da fonte}{will probably never be entirely factual and completely devoid of bias}{Broader Impact.}
\end{itemize}

\campo{Limitações declaradas e análise crítica}
Autores: fontes externas, inclusive a Wikipédia, podem conter inexatidões e vieses (Broader Impact). \textbf{Análise da ficha:} a arquitetura treinada do artigo não é idêntica a um pipeline que apenas inclui trechos em prompts de um LLM; os ganhos não são transferíveis automaticamente.

\campo{Aproveitamento na dissertação · proposta}
Referencial teórico: definir RAG e distinguir a formulação original da implementação usada na dissertação. Metodologia: justificar medir geração e recuperação sem presumir eliminação de alucinações.

\registro{Fonte consultada nesta preparação: PDF publicado; consulta focal de Resumo, seções 2--4, Discussão e Broader Impact. Apêndices não fichados. Verificação estrutural automática indisponível; usados localizadores textuais por seção. Excertos em inglês. Chave: \texttt{lewisEtAl2020rag}.}
